% !TeX root = ../../Book2.tex
\section{Gauss 函数与自对偶性}
\label{b2:sec:0602}

Poisson 求和把一个函数的格点和转成其变换的格点和。Gauss 函数的变换仍为 Gauss 函数，因此可以显式比较大尺度和小尺度，并把这种比较放回热传播。

\subsection{Gauss 函数的 Fourier 变换}
\label{b2:subsec:060201}

记 \(g_t(x)=e^{-\pi t x^2}\)，其中 \(t>0\)。由\cref{b2:thm:03-gaussian}及伸缩公式，
\[
 \widehat {g_t}(\xi)=t^{-1/2}e^{-\pi\xi^2/t}.
\]
与保持单位质量的 \(G_\varepsilon\) 相比，\(g_t\) 在原点的高度固定为 \(1\)，总积分为 \(t^{-1/2}\)。将两类参数混用会漏掉平方根因子；这里分别用 \(t\) 和 \(\varepsilon\) 保留区别。

\ProseParagraphBreak

平移和调制给出更一般的波包公式。对实数 \(a,b\)，
\[
 \mathcal F\bigl(e^{-\pi t(x-a)^2}e^{2\pi\mathrm{i}bx}\bigr)(\xi)
 =t^{-1/2}e^{-2\pi\mathrm{i}a(\xi-b)}e^{-\pi(\xi-b)^2/t}.
\]
因此空间中心 \(a\) 产生相位，调制频率 \(b\) 产生频率中心，而宽度参数倒置。每一项都来自\chapref{b2:ch:03}已经证明的运算规则。

\subsection{自对偶结构}
\label{b2:subsec:060202}

对本卷 Fourier 变换已定义的函数或分布 \(f\)，若 \(\widehat f=f\) 在同一空间中成立，称 \(f\) 对当前 Fourier 变换\term[zi-duiou]{自对偶}[self-dual]。Gauss 函数 \(g_1\) 满足这一条件。一般 \(g_t\) 的参数变为 \(1/t\)，并附带 \(t^{-1/2}\)；只有 \(t=1\) 保持同一个归一化函数。

\ProseParagraphBreak

\subsecref{b2:subsec:050204}的 Hermite 函数提供了正 Gauss 函数之外的自对偶例子。\cref{b2:thm:05-hermite}给出 \(\widehat h_n=(-\mathrm i)^nh_n\)，故对每个 \(j\in\mathbb N_0\)，\(h_{4j}\) 都自对偶，有限复线性组合仍然自对偶。相反，\(h_{4j+2}\) 变换后变号，奇数指标对应虚特征值。自对偶是一项线性变换关系，不意味着非负性或概率密度性质。

\subsection{热核解释}
\label{b2:subsec:060203}

\chapref{b2:ch:05}已使用 \(H_t(x)=(4\pi t)^{-1/2}e^{-x^2/(4t)}\)，其变换为 \(e^{-4\pi^2t\xi^2}\)。在实线上周期化，得到
\begin{equation}\label{b2:eq:06-periodic-heat}
 K_t(x)\coloneqq\sum_{n\in\mathbb Z}H_t(x+n)
 =\sum_{k\in\mathbb Z}e^{-4\pi^2tk^2}e^{2\pi\mathrm{i}kx}.
\end{equation}
第二个表达由 Poisson 求和得到；对固定 \(t>0\)，两种表达都绝对收敛。函数 \(K_t\) 称为圆周上的\term[rehe]{热核}[heat kernel]。
\symidx{\(K_t\)：圆周热核}

\begin{proposition}\label{b2:prop:06-heat-kernel}
对 \(t,s>0\)，令 \(K_t(x)=\sum_{n\in\mathbb Z}(4\pi t)^{-1/2}e^{-(x+n)^2/(4t)}\)，在 \(\mathbb T=\mathbb R/\mathbb Z\) 上取总质量为 \(1\) 的 Lebesgue 测度，\(*\) 表示圆周卷积。圆周热核 \(K_t\) 非负、积分为 \(1\)，满足 \(K_t*K_s=K_{t+s}\) 和 \(\partial_tK_t=\partial_x^2K_t\)。对 \(f\in L^p(\mathbb T;\mathbb C)\)、\(1\leq p<\infty\)，当 \(t\downarrow0\) 时有 \(K_t*f\to f\) 于 \(L^p\)；对 \(f\in C(\mathbb T;\mathbb C)\) 则一致收敛。
\end{proposition}
\begin{proof}
非负性由实线表达给出，非负项的 Tonelli 换序将其圆周积分化为实线 Gauss 积分，故积分为 \(1\)。两个核均可积，圆周卷积的 Fourier 系数等于两核系数之积，即 \(e^{-4\pi^2(t+s)k^2}\)；由\cref{b2:cor:02-uniqueness}的周期系数唯一性及两侧连续性得到半群关系。对 \(t\geq\varepsilon>0\)，任意阶 \(x,t\) 导数只在系数上增加一个固定次幂的 \(k\)，仍被可和的指数尾统一控制，故可逐项求导，并得到热方程。

为证初值逼近，对 \(0<\delta<1/2\)，圆周上离原点至少为 \(\delta\) 的区域在周期化前对应到实线中 \(|y|\geq\delta\) 的子集。因此
\[
 \int_{\mathrm{d}_{\mathbb T}(x,0)\geq\delta}K_t(x)\,\mathrm{d}x
 \leq\int_{|y|\geq\delta}H_t(y)\,\mathrm{d}y\longrightarrow0.
\]
所以它满足\cref{b2:prop:02-positive-approximation}中的三项条件：非负、单位质量以及远离原点的质量趋零。该命题对任意趋零的正时间列都适用，因而得到 \(t\downarrow0\) 时的全部收敛结论。
\end{proof}

小时间时，实线表达用少量邻近格点就能描述局部尖峰；大时间时，频率表达只需保留少量低频，就能看出函数趋向常数 \(1\)。两个公式描述同一个核，却在相反的尺度下各自简洁。下一节的 \(\theta\) 变换正是这种尺度对换在 \(x=0\) 处的表达。

\ProseParagraphBreak

这种比较可以给出两端都可核算的误差。对 \(t\geq t_0>0\)，由 \(k^2\geq1+3(k-1)\)、\(k\geq1\)，频率表达给出
\[
 \|K_t-1\|_\infty
 \leq2\sum_{k=1}^\infty e^{-4\pi^2tk^2}
 \leq\frac{2e^{-4\pi^2t}}{1-e^{-12\pi^2t}}.
\]
因此向均匀分布趋近的首个指数由 \(k=\pm1\) 决定。对 \(f\in L^2(\mathbb T;\mathbb C)\)，记平均值 \(\overline f=\int_0^1f\)，这里上横线表示平均而非复共轭。\cref{b2:cor:02-parseval}将 \(L^2\) 范数平方写成全部 Fourier 系数模平方之和，给出精确的能量表达
\[
 \|K_t*f-\overline f\|_2^2
 =\sum_{k\ne0}e^{-8\pi^2tk^2}|\widehat f(k)|^2
 \leq e^{-8\pi^2t}\|f-\overline f\|_2^2.
\]
取 \(f(x)=e^{2\pi\mathrm{i}x}\) 时达到等号，所以能量范数中的衰减率 \(4\pi^2\) 不能统一增大。常数模态保留平均，其余模态耗散。

\ProseParagraphBreak

对于小时间，固定 \(0<a<1/2\) 并取 \(|x|\leq a\)。零格点项为 \(H_t(x)\)，其余项满足 \(|x+n|\geq|n|-a\)。对 \(n\geq1\)，
\[
 (n-a)^2\geq(1-a)^2+(3-2a)(n-1),
\]
故
\[
 0\leq K_t(x)-H_t(x)
 \leq\frac{2(4\pi t)^{-1/2}e^{-(1-a)^2/(4t)}}
              {1-e^{-(3-2a)/(4t)}}.
\]
当 \(0<t\leq1\) 时，分母有依赖 \(a\) 的统一正下界。于是邻近原点的圆周传播，与实线 Gauss 传播之间只有指数小的误差。若接近半周期端点，就可能有两个相邻格点同样重要；单独保留零格点的近似范围必须随之改变。
